% Generated draft: requires mathematical and source review.
\documentclass[11pt]{article}
\usepackage[margin=25mm]{geometry}
\usepackage{fontspec}
\setmainfont{Latin Modern Roman}
\usepackage{amsmath,amssymb,mathtools,unicode-math}
\setmathfont{Latin Modern Math}
\usepackage{enumitem,needspace,xurl,hyperref}
\setlength{\parindent}{0pt}
\setlength{\parskip}{5pt}
\setlength{\emergencystretch}{3em}
\newcommand{\N}{\mathbb{N}}
\newcommand{\Z}{\mathbb{Z}}
\newcommand{\Q}{\mathbb{Q}}
\newcommand{\R}{\mathbb{R}}
\newcommand{\C}{\mathbb{C}}
\newcommand{\F}{\mathbb{F}}
\newcommand{\A}{\mathbb{A}}
\newcommand{\PP}{\mathbb{P}}
\newcommand{\E}{\mathbb{E}}
\newcommand{\eps}{\varepsilon}
\newcommand{\dd}{\,\mathrm{d}}
\begin{document}
\section*{UnsolvedMath ID: 20001938}
\textbf{Source problem code:} catalogue entry
\subsection{Definitions and mathematical statement}
Let $M^n$, $n\ge2$, be a smooth manifold. A projective structure $[\nabla]$ is an equivalence class of torsion-free affine connections with the same unparametrized geodesics; representatives satisfy
\[
 \widehat\nabla_XY=\nabla_XY+\Upsilon(X)Y+\Upsilon(Y)X
\]
for some one-form $\Upsilon$. A compatible metric is a nondegenerate pseudo-Riemannian metric whose Levi-Civita connection belongs to $[\nabla]$. Positive definiteness is an additional condition, not implied by nondegeneracy.

\textbf{Problem 12.0.5.} Construct projectively invariant scalar differential tests whose simultaneous vanishing characterizes existence of a compatible metric locally in a neighborhood of almost every point. The source does not specify whether ``almost every'' is measure theoretic or means a generic regular locus, nor whether its metric is required to be positive definite. A frame-dependent collection of minors with invariant common zero set is not automatically a collection of scalar-valued projective invariants.

\textbf{Precise form of the supplied regular-stratum reduction.} The weighted projective metrizability equation may be written
\[
 \nabla_a\sigma^{bc}-\frac1{n+1}
 \bigl(\delta_a^b\nabla_d\sigma^{dc}
       +\delta_a^c\nabla_d\sigma^{bd}\bigr)=0,
 \qquad \sigma\in\Gamma(S^2TM\otimes\mathcal E(-2)),
\]
where $\mathcal E(-2)$ is the standard projective density bundle of weight $-2$. Nondegenerate solutions correspond locally to compatible pseudo-Riemannian metrics. The aid assumes the \emph{exact metrizability prolongation connection} $D$ on a bundle $\mathcal V$ of rank $N=(n+1)(n+2)/2$, whose parallel sections correspond to solutions; this is not silently replaced by an uncorrected normal tractor connection. Let $\pi:\mathcal V\to S^2TM\otimes\mathcal E(-2)$ be its top-slot projection and $R^D$ its curvature. Choose an auxiliary affine connection $B$ on $TM$ and use $D$ and $B$ to form full covariant derivatives $\mathcal D^jR^D$. Define the cumulative bundle maps
\[
 H_j(s)=\bigl(R^D s,(\mathcal D R^D)s,\ldots,
                         (\mathcal D^jR^D)s\bigr),
 \qquad K_j=\ker H_j,\quad 0\le j\le N-1.
\]
All tensor components and all derivative directions are retained. On a sufficiently small simply connected neighborhood where all the ranks $\operatorname{rank}H_0,\ldots,\operatorname{rank}H_{N-1}$ are constant, the aid asserts that $K=K_{N-1}$ is the stabilized flat parallel-section bundle. Consequently all $N\times N$ minors of a matrix for $H_{N-1}$ vanish exactly when a nonzero local metrizability solution exists.

This nonzero-solution test is weaker than metric existence. At $x$, define the density-valued polynomial on $K_x$ by
\[
 P_x(s)=\det(\pi_xs).
\]
The aid asserts that a compatible pseudo-Riemannian metric exists near $x$ precisely when $P_x$ is not the zero polynomial, equivalently some $s\in K_x$ has nondegenerate top slot. A Riemannian conclusion additionally requires a definite top slot. Thus vanishing of the curvature minors alone does not solve the metric problem.

Changing $B$ is claimed to transform cumulative maps by invertible block lower-triangular maps with identity diagonal; changing projective representative additionally changes the splitting invertibly. The invariant object asserted by the aid is therefore the kernel filtration and the locally generated determinantal ideal, not each numerical minor. The common constant-rank locus is stated to be open dense; rank-changing strata are unresolved. In the smooth category, open dense need not have full measure, so the original ``almost every point'' qualification is not silently settled.

\textbf{Remaining source context.} The source reports a dimension-two solution [Br-Du-Eas2009], describing its simplest invariant as degree five in derivatives, and suggests that stronger overdeterminacy in higher dimensions may simplify early obstructions [Dun-Eas2014]. ``Degree five'' is retained without deciding whether it denotes differential order or polynomial degree. It reports generic uniqueness up to constant scaling [Matveev2012a], classical nonproportional examples due to Lagrange, Beltrami and Dini, local Riemannian classification [Levi-Civita1896], and arbitrary-signature local classification [Bol-Mat2013]. Its further broad problem is to understand which closed manifolds admit nonproportional projectively equivalent metrics of indefinite signature. The final phrase introducing a ``simplest version'' has no ensuing question in this extract; no topology conjecture is supplied in its place. Historical assertions of how little is known globally are source descriptions, not a current assessment.
\subsection{Short English statement}
The problem asks for invariant scalar obstructions that exactly detect local metric realizations of a projective structure. The supplied aid gives a curvature-jet kernel test on constant-rank strata, but a nonzero solution must also contain a nondegenerate metric component. Rank-changing points, the meaning of almost every point, and conversion from determinantal ideals to scalar invariants remain qualifications.
\subsection{Sources}
\begin{itemize}
\item Local UnsolvedMath registry, permanent ID 20001938, source code catalogue entry; canonical \texttt{aim-geometry-notes.json}, notes[275].
\item AIM \emph{Open Problems and Questions about Geodesics}, Problem 12.0.5 and following discussion: \url{https://aimath.org/pastworkshops/geodesicsproblems.pdf}.
\item Supplied labels [Br-Du-Eas2009], [Dun-Eas2014], [Matveev2012a], [Levi-Civita1896], [Bol-Mat2013]. Full bibliographic entries are absent. Historical examples are attributed to Lagrange, Beltrami and Dini.
\item Registry research attempt 1, dated 2026-08-14, partial result, rigor ``proved'', novelty confidence low. Its explicitly machine-generated, non-peer-reviewed research aid requires independent verification before citation.
\end{itemize}
\end{document}
